\section{引言}

量子色动力学（QCD）中最轻的束缚态——$\pi$介子——扮演着基本性的角色，因为它是动力学手征对称性破缺（DCSB）的南部–戈德斯通玻色子。
对$\pi$介子与K介子结构的研究揭示了DCSB的物理，有助于厘清DCSB与由夸克质量引起的手征对称性破缺二者的相对影响，并对理解非微扰QCD十分重要。
研究$\pi$介子部分子分布函数（PDF）对刻画$\pi$介子结构、进一步理解DCSB与非微扰QCD都很重要。
目前我们对$\pi$介子PDF的认识仍少于对核子PDF的认识，原因在于实验数据集较少，尤其是海夸克与胶子分布方面。
未来建于布鲁克海文国家实验室的美国电子离子对撞机（EIC）~\cite{Accardi:2012qut}将增进我们对$\pi$介子结构的认识~\cite{Arrington:2021biu,Aguilar:2019teb}。
%
在中国，类似的装置——中国电子离子对撞机（EicC）~\cite{Anderle:2021wcy}——也已规划，将对$\pi$介子胶子与海夸克分布的研究产生影响。
%
在欧洲，COMPASS++/AMBER 的 Drell--Yan 与$J/\psi$产生实验~\cite{Denisov:2018unj}旨在同时增进我们对$\pi$介子胶子与夸克PDF的认识。

$\pi$介子PDF的全局分析大多依赖 Drell--Yan 数据。
早期$\pi$介子PDF研究主要基于$\pi$介子诱发的 Drell--Yan 数据，并利用$J/\psi$产生或直接光子产生数据来约束$\pi$介子胶子PDF~\cite{Owens:1984zj,Aurenche:1989sx,Sutton:1991ay,Gluck:1991ey,Gluck:1999xe}。
也有较新的工作，例如 Bourrely 与 Soffer~\cite{Bourrely:2018yck} 基于 Drell--Yan $\pi^+W$ 数据提取$\pi$介子PDF。
JAM 合作组~\cite{Barry:2018ort,Cao:2021aci}用蒙特卡罗方法分析 Drell--Yan $\pi A$ 及来自 HERA 的领头中子电产生数据以进入较低-$x$区间，并发现胶子在$\pi$介子中所携带的动量份额（约$30\%$）显著高于仅凭 Drell--Yan 数据的推断。
xFitter 组~\cite{Novikov:2020snp}利用其开源 QCD 拟合框架分析 Drell--Yan $\pi A$ 与光产生数据，发现这些数据能很好地约束价分布，但对海夸克与胶子分布的约束尚不足以精确确定它们。
文献~\cite{Chang:2020rdy}的分析表明，$\pi$介子诱发的$J/\psi$产生数据为$\pi$介子PDF提供了额外约束，尤其是大-$x$区的$\pi$介子胶子PDF。
总而言之，来自实验数据全局分析的$\pi$介子价夸克分布比胶子分布受到更好的约束。
在等待更多实验数据集的同时，从理论一侧开展$\pi$介子胶子分布研究可为实验提供有用的信息。

$\pi$介子胶子PDF很少被连续QCD唯象模型或格点QCD（LQCD）模拟所研究。
大多数模型研究只预言$\pi$介子的价夸克分布~\cite{Nam:2012vm,Watanabe:2016lto,Watanabe:2017pvl,Hutauruk:2016sug,Lan:2019vui,Lan:2019rba,deTeramond:2018ecg,Watanabe:2019zny,Han:2018wsw,Chang:2014lva,Chang:2014gga,Chen:2016sno,Shi:2018mcb,Bednar:2018mtf,Ding:2019lwe}，
而胶子与海PDF则由 Dyson--Schwinger 方程（DSE）连续方法预言~\cite{Freese:2021zne}。基于彩虹梯截断实现的 DSE 对$\pi$介子胶子PDF的预言，与 JAM 的$\pi$介子胶子PDF结果~\cite{Barry:2018ort,Cao:2021aci}在两个标准差内一致。
LQCD 提供第一性原理计算来改进我们对非微扰$\pi$介子胶子结构的认识；
然而迄今只有两项工作确定过$\pi$介子胶子PDF的一阶矩~\cite{Meyer:2007tm,Shanahan:2018pib}。
2000年的一项早期计算采用 quenched QCD、Wilson 费米子作用量、晶格间距$a=0.093$~fm、格点体积$L^3\times T=24^4$、高达890~MeV的$\pi$介子质量以及3066个组型，给出$\mu^2=4\text{ GeV}^2$处$\langle x \rangle_g=0.37(8)(12)$~\cite{Meyer:2007tm}。
2018年的一项较新研究采用$N_f=2+1$ clover 费米子作用量、晶格间距$a=0.1167(16)$~fm、更大的格点$32^3\times 96$、450~MeV的$\pi$介子质量以及572{,}663次测量，给出了更大的一阶矩结果：$\mu^2=4\text{ GeV}^2$处$\langle x \rangle_g=0.61(9)$~\cite{Shanahan:2018pib}。
%
原则上，一系列矩可用于重构PDF。
尽管已有$\pi$介子胶子PDF一阶矩的计算，但要完成这样的重构，获得足够多高阶矩的希望渺茫。
因此需要对$\pi$介子胶子PDF的$x$依赖做直接的格点计算。

近年来，随着大动量有效理论（LaMET）~\cite{Ji:2013dva,Ji:2014gla,Ji:2017rah}的提出，格点QCD中$x$依赖强子结构的计算日益增多。
LaMET 方法计算格点准分布函数——定义为等时空间间隔算符的矩阵元——然后取无穷大动量极限来提取光锥分布。
准PDF可通过因子化定理与$P_z$无关的光锥PDF相联系，从中分解出微扰匹配系数以及受强子动量压低的修正~\cite{Ji:2014gla}。
该因子化可在微扰论中精确计算~\cite{Ma:2017pxb,Liu:2019urm}。
基于准PDF方法，已有许多针对核子与介子PDF以及广义部分子分布（GPD）的格点工作~\cite{Lin:2013yra,Lin:2014zya,Chen:2016utp,Lin:2017ani,Alexandrou:2015rja,Alexandrou:2016jqi,Alexandrou:2017huk,Chen:2017mzz,Alexandrou:2018pbm,Chen:2018xof,Chen:2018fwa,Alexandrou:2018eet,Lin:2018qky,Fan:2018dxu,Liu:2018hxv,Wang:2019tgg,Lin:2019ocg,Chen:2019lcm,Lin:2020reh,Chai:2020nxw,Bhattacharya:2020cen,Lin:2020ssv,Zhang:2020dkn,Li:2020xml,Fan:2020nzz,Gao:2020ito,Lin:2020fsj,Zhang:2020rsx,Alexandrou:2020qtt,Alexandrou:2020zbe,Lin:2020rxa,Gao:2021hxl,Lin:2020rut}。
在格点QCD中获取光锥PDF的其他途径还有``好格点截面''~\cite{Ma:2017pxb,Bali:2017gfr,Bali:2018spj,Sufian:2019bol,Sufian:2020vzb}与赝PDF方法~\cite{Orginos:2017kos,Karpie:2017bzm,Karpie:2018zaz,Karpie:2019eiq,Joo:2019jct,Joo:2019bzr,Radyushkin:2018cvn,Zhang:2018ggy,Izubuchi:2018srq,Joo:2020spy,Bhat:2020ktg,Fan:2020cpa,Sufian:2020wcv,Karthik:2021qwz}。
然而，关于$\pi$介子PDF的$x$依赖，LQCD 计算此前只做过价夸克分布~\cite{Chen:2018fwa,Sufian:2019bol,Izubuchi:2019lyk,Joo:2019bzr,Sufian:2020vzb,Shugert:2020tgq,Gao:2020ito}。

直到最近，当赝PDF~\cite{Balitsky:2019krf}与准PDF~\cite{Zhang:2018diq,Wang:2019tgg}途径所需的一圈匹配关系被算出后，胶子PDF的格点计算才成为可能。
两种途径都使直接计算$\pi$介子胶子PDF的$x$依赖变得可行。
本工作中，我们应用赝PDF方法并采用比值重正化方案，以回避计算胶子重正化因子的困难。
利用赝PDF方法从Ioffe时间分布（ITD）获取光锥PDF已有一套成熟的流程：经演化与方案转换两步进行匹配~\cite{Radyushkin:2018cvn}。
使用赝PDF方法还允许我们利用小Ioffe时间处所有提升动量的格点关联函数。
此前已有若干成功的核子同位旋矢量
PDF~\cite{Orginos:2017kos,Joo:2019jct,Joo:2020spy,Bhat:2020ktg}与$\pi$介子价夸克PDF~\cite{Joo:2019bzr}的赝PDF计算。
最早的计算在 quenched 格点上完成~\cite{Orginos:2017kos}，随后$\pi$介子质量逐步取到更接近物理值~\cite{Joo:2019jct,Joo:2019bzr,Joo:2020spy}，物理$\pi$介子质量的计算也于最近实现~\cite{Bhat:2020ktg}。
文献~\cite{Joo:2019jct,Joo:2019bzr,Joo:2020spy,Bhat:2020ktg}的格点计算PDF与全局分析PDF符合良好。

本工作中，我们用赝PDF方法首次完整计算了$\pi$介子胶子分布的$x$依赖，使用了0.12与0.15~fm两种晶格间距和690、310、220~MeV三种$\pi$介子质量。
本文其余部分安排如下。
第~\ref{sec:cal-details}节介绍由约化ITD获取光锥胶子PDF的流程、格点模拟的数值设置，以及如何从格点计算的关联函数中提取约化ITD。
第~\ref{sec:results}节把我们从格点计算最终确定的$\pi$介子胶子PDF与 NLO 的 xFitter~\cite{Novikov:2020snp}和 JAM $\pi$介子胶子PDF~\cite{Barry:2018ort,Cao:2021aci}进行比较；
文中研究了各步骤引入的系统误差，并考察了晶格间距与$\pi$介子质量依赖。
